\IfFileExists{stacks-project.cls}{%
\documentclass{stacks-project}
}{%
\documentclass{amsart}
}

% For dealing with references we use the comment environment
\usepackage{verbatim}
\newenvironment{reference}{\comment}{\endcomment}
%\newenvironment{reference}{}{}
\newenvironment{slogan}{\comment}{\endcomment}
\newenvironment{history}{\comment}{\endcomment}

% For commutative diagrams we use Xy-pic
\usepackage[all]{xy}

% We use 2cell for 2-commutative diagrams.
\xyoption{2cell}
\UseAllTwocells

% We use multicol for the list of chapters between chapters
\usepackage{multicol}

% This is generally recommended for better output
\usepackage{lmodern}
\usepackage[T1]{fontenc}

% For cross-file-references
\usepackage{xr-hyper}

% Package for hypertext links:
\usepackage{hyperref}

% For any local file, say "hello.tex" you want to link to please
% use \externaldocument[hello-]{hello}
\externaldocument[introduction-]{introduction}
\externaldocument[conventions-]{conventions}
\externaldocument[sets-]{sets}
\externaldocument[categories-]{categories}
\externaldocument[topology-]{topology}
\externaldocument[sheaves-]{sheaves}
\externaldocument[sites-]{sites}
\externaldocument[stacks-]{stacks}
\externaldocument[fields-]{fields}
\externaldocument[algebra-]{algebra}
\externaldocument[brauer-]{brauer}
\externaldocument[homology-]{homology}
\externaldocument[derived-]{derived}
\externaldocument[simplicial-]{simplicial}
\externaldocument[more-algebra-]{more-algebra}
\externaldocument[smoothing-]{smoothing}
\externaldocument[modules-]{modules}
\externaldocument[sites-modules-]{sites-modules}
\externaldocument[injectives-]{injectives}
\externaldocument[cohomology-]{cohomology}
\externaldocument[sites-cohomology-]{sites-cohomology}
\externaldocument[dga-]{dga}
\externaldocument[dpa-]{dpa}
\externaldocument[sdga-]{sdga}
\externaldocument[hypercovering-]{hypercovering}
\externaldocument[schemes-]{schemes}
\externaldocument[constructions-]{constructions}
\externaldocument[properties-]{properties}
\externaldocument[morphisms-]{morphisms}
\externaldocument[coherent-]{coherent}
\externaldocument[divisors-]{divisors}
\externaldocument[limits-]{limits}
\externaldocument[varieties-]{varieties}
\externaldocument[topologies-]{topologies}
\externaldocument[descent-]{descent}
\externaldocument[perfect-]{perfect}
\externaldocument[more-morphisms-]{more-morphisms}
\externaldocument[flat-]{flat}
\externaldocument[groupoids-]{groupoids}
\externaldocument[more-groupoids-]{more-groupoids}
\externaldocument[etale-]{etale}
\externaldocument[chow-]{chow}
\externaldocument[intersection-]{intersection}
\externaldocument[pic-]{pic}
\externaldocument[weil-]{weil}
\externaldocument[adequate-]{adequate}
\externaldocument[dualizing-]{dualizing}
\externaldocument[duality-]{duality}
\externaldocument[discriminant-]{discriminant}
\externaldocument[derham-]{derham}
\externaldocument[local-cohomology-]{local-cohomology}
\externaldocument[algebraization-]{algebraization}
\externaldocument[curves-]{curves}
\externaldocument[resolve-]{resolve}
\externaldocument[models-]{models}
\externaldocument[functors-]{functors}
\externaldocument[equiv-]{equiv}
\externaldocument[pione-]{pione}
\externaldocument[etale-cohomology-]{etale-cohomology}
\externaldocument[proetale-]{proetale}
\externaldocument[relative-cycles-]{relative-cycles}
\externaldocument[more-etale-]{more-etale}
\externaldocument[trace-]{trace}
\externaldocument[crystalline-]{crystalline}
\externaldocument[spaces-]{spaces}
\externaldocument[spaces-properties-]{spaces-properties}
\externaldocument[spaces-morphisms-]{spaces-morphisms}
\externaldocument[decent-spaces-]{decent-spaces}
\externaldocument[spaces-cohomology-]{spaces-cohomology}
\externaldocument[spaces-limits-]{spaces-limits}
\externaldocument[spaces-divisors-]{spaces-divisors}
\externaldocument[spaces-over-fields-]{spaces-over-fields}
\externaldocument[spaces-topologies-]{spaces-topologies}
\externaldocument[spaces-descent-]{spaces-descent}
\externaldocument[spaces-perfect-]{spaces-perfect}
\externaldocument[spaces-more-morphisms-]{spaces-more-morphisms}
\externaldocument[spaces-flat-]{spaces-flat}
\externaldocument[spaces-groupoids-]{spaces-groupoids}
\externaldocument[spaces-more-groupoids-]{spaces-more-groupoids}
\externaldocument[bootstrap-]{bootstrap}
\externaldocument[spaces-pushouts-]{spaces-pushouts}
\externaldocument[spaces-chow-]{spaces-chow}
\externaldocument[groupoids-quotients-]{groupoids-quotients}
\externaldocument[spaces-more-cohomology-]{spaces-more-cohomology}
\externaldocument[spaces-simplicial-]{spaces-simplicial}
\externaldocument[spaces-duality-]{spaces-duality}
\externaldocument[formal-spaces-]{formal-spaces}
\externaldocument[restricted-]{restricted}
\externaldocument[spaces-resolve-]{spaces-resolve}
\externaldocument[formal-defos-]{formal-defos}
\externaldocument[defos-]{defos}
\externaldocument[cotangent-]{cotangent}
\externaldocument[examples-defos-]{examples-defos}
\externaldocument[algebraic-]{algebraic}
\externaldocument[examples-stacks-]{examples-stacks}
\externaldocument[stacks-sheaves-]{stacks-sheaves}
\externaldocument[criteria-]{criteria}
\externaldocument[artin-]{artin}
\externaldocument[quot-]{quot}
\externaldocument[stacks-properties-]{stacks-properties}
\externaldocument[stacks-morphisms-]{stacks-morphisms}
\externaldocument[stacks-limits-]{stacks-limits}
\externaldocument[stacks-cohomology-]{stacks-cohomology}
\externaldocument[stacks-perfect-]{stacks-perfect}
\externaldocument[stacks-introduction-]{stacks-introduction}
\externaldocument[stacks-more-morphisms-]{stacks-more-morphisms}
\externaldocument[stacks-geometry-]{stacks-geometry}
\externaldocument[moduli-]{moduli}
\externaldocument[moduli-curves-]{moduli-curves}
\externaldocument[examples-]{examples}
\externaldocument[exercises-]{exercises}
\externaldocument[guide-]{guide}
\externaldocument[desirables-]{desirables}
\externaldocument[coding-]{coding}
\externaldocument[obsolete-]{obsolete}
\externaldocument[fdl-]{fdl}
\externaldocument[index-]{index}

% Theorem environments.
%
\theoremstyle{plain}
\newtheorem{theorem}[subsection]{Theorem}
\newtheorem{proposition}[subsection]{Proposition}
\newtheorem{lemma}[subsection]{Lemma}

\theoremstyle{definition}
\newtheorem{definition}[subsection]{Definition}
\newtheorem{example}[subsection]{Example}
\newtheorem{exercise}[subsection]{Exercise}
\newtheorem{situation}[subsection]{Situation}

\theoremstyle{remark}
\newtheorem{remark}[subsection]{Remark}
\newtheorem{remarks}[subsection]{Remarks}

\numberwithin{equation}{subsection}

% Macros
%
\def\lim{\mathop{\mathrm{lim}}\nolimits}
\def\colim{\mathop{\mathrm{colim}}\nolimits}
\def\Spec{\mathop{\mathrm{Spec}}}
\def\Hom{\mathop{\mathrm{Hom}}\nolimits}
\def\Ext{\mathop{\mathrm{Ext}}\nolimits}
\def\SheafHom{\mathop{\mathcal{H}\!\mathit{om}}\nolimits}
\def\SheafExt{\mathop{\mathcal{E}\!\mathit{xt}}\nolimits}
\def\Sch{\mathit{Sch}}
\def\Mor{\mathop{\mathrm{Mor}}\nolimits}
\def\Ob{\mathop{\mathrm{Ob}}\nolimits}
\def\Sh{\mathop{\mathit{Sh}}\nolimits}
\def\NL{\mathop{N\!L}\nolimits}
\def\CH{\mathop{\mathrm{CH}}\nolimits}
\def\proetale{{pro\text{-}\acute{e}tale}}
\def\etale{{\acute{e}tale}}
\def\QCoh{\mathit{QCoh}}
\def\Ker{\mathop{\mathrm{Ker}}}
\def\Im{\mathop{\mathrm{Im}}}
\def\Coker{\mathop{\mathrm{Coker}}}
\def\Coim{\mathop{\mathrm{Coim}}}

% Boxtimes
%
\DeclareMathSymbol{\boxtimes}{\mathbin}{AMSa}{"02}

%
% Macros for moduli stacks/spaces
%
\def\QCohstack{\mathcal{QC}\!\mathit{oh}}
\def\Cohstack{\mathcal{C}\!\mathit{oh}}
\def\Spacesstack{\mathcal{S}\!\mathit{paces}}
\def\Quotfunctor{\mathrm{Quot}}
\def\Hilbfunctor{\mathrm{Hilb}}
\def\Curvesstack{\mathcal{C}\!\mathit{urves}}
\def\Polarizedstack{\mathcal{P}\!\mathit{olarized}}
\def\Complexesstack{\mathcal{C}\!\mathit{omplexes}}
% \Pic is the operator that assigns to X its picard group, usage \Pic(X)
% \Picardstack_{X/B} denotes the Picard stack of X over B
% \Picardfunctor_{X/B} denotes the Picard functor of X over B
\def\Pic{\mathop{\mathrm{Pic}}\nolimits}
\def\Picardstack{\mathcal{P}\!\mathit{ic}}
\def\Picardfunctor{\mathrm{Pic}}
\def\Deformationcategory{\mathcal{D}\!\mathit{ef}}
\usepackage{cleveref}
\crefname{originaldefinition}{original definition}{original definitions}
\Crefname{originaldefinition}{Original definition}{Original definitions}
\makeatletter
\newlabel{example-reduction-affine-formal-spectrum}{{0GB6}{source}{}{sourcepointer-1}{}}
\newlabel{example-reduction-affine-formal-spectrum@cref}{{[originaldefinition][0][]0GB6}{[1][0][]source}{}{}{}}
\makeatother
\crefname{originallemma}{original lemma}{original lemmas}
\Crefname{originallemma}{Original lemma}{Original lemmas}
\makeatletter
\newlabel{lemma-completed-tensor-product}{{0GB4}{source}{}{sourcepointer-2}{}}
\newlabel{lemma-completed-tensor-product@cref}{{[originallemma][0][]0GB4}{[1][0][]source}{}{}{}}
\makeatother
\crefname{originallemma}{original lemma}{original lemmas}
\Crefname{originallemma}{Original lemma}{Original lemmas}
\makeatletter
\newlabel{lemma-completed-tensor-product-weakly-adic}{{0GXJ}{source}{}{sourcepointer-3}{}}
\newlabel{lemma-completed-tensor-product-weakly-adic@cref}{{[originallemma][0][]0GXJ}{[1][0][]source}{}{}{}}
\makeatother
\crefname{originallemma}{original lemma}{original lemmas}
\Crefname{originallemma}{Original lemma}{Original lemmas}
\makeatletter
\newlabel{lemma-countable-affine-formal-algebraic-space}{{0AII}{source}{}{sourcepointer-4}{}}
\newlabel{lemma-countable-affine-formal-algebraic-space@cref}{{[originallemma][0][]0AII}{[1][0][]source}{}{}{}}
\makeatother
\crefname{originallemma}{original lemma}{original lemmas}
\Crefname{originallemma}{Original lemma}{Original lemmas}
\makeatletter
\newlabel{lemma-countably-indexed}{{0AN1}{source}{}{sourcepointer-5}{}}
\newlabel{lemma-countably-indexed@cref}{{[originallemma][0][]0AN1}{[1][0][]source}{}{}{}}
\makeatother
\crefname{originallemma}{original lemma}{original lemmas}
\Crefname{originallemma}{Original lemma}{Original lemmas}
\makeatletter
\newlabel{lemma-factor-through-thickening}{{0AIA}{source}{}{sourcepointer-6}{}}
\newlabel{lemma-factor-through-thickening@cref}{{[originallemma][0][]0AIA}{[1][0][]source}{}{}{}}
\makeatother
\crefname{originallemma}{original lemma}{original lemmas}
\Crefname{originallemma}{Original lemma}{Original lemmas}
\makeatletter
\newlabel{lemma-fibre-product-affines-over-separated}{{0AN3}{source}{}{sourcepointer-7}{}}
\newlabel{lemma-fibre-product-affines-over-separated@cref}{{[originallemma][0][]0AN3}{[1][0][]source}{}{}{}}
\makeatother
\crefname{originallemma}{original lemma}{original lemmas}
\Crefname{originallemma}{Original lemma}{Original lemmas}
\makeatletter
\newlabel{lemma-morphism-between-formal-spectra}{{0AN0}{source}{}{sourcepointer-8}{}}
\newlabel{lemma-morphism-between-formal-spectra@cref}{{[originallemma][0][]0AN0}{[1][0][]source}{}{}{}}
\makeatother
\crefname{originallemma}{original lemma}{original lemmas}
\Crefname{originallemma}{Original lemma}{Original lemmas}
\makeatletter
\newlabel{lemma-reduction-smooth}{{0GB7}{source}{}{sourcepointer-9}{}}
\newlabel{lemma-reduction-smooth@cref}{{[originallemma][0][]0GB7}{[1][0][]source}{}{}{}}
\makeatother
\crefname{originallemma}{original lemma}{original lemmas}
\Crefname{originallemma}{Original lemma}{Original lemmas}
\makeatletter
\newlabel{spaces-morphisms-lemma-finite-type-local}{{040Y}{source}{}{sourcepointer-10}{}}
\newlabel{spaces-morphisms-lemma-finite-type-local@cref}{{[originallemma][0][]040Y}{[1][0][]source}{}{}{}}
\makeatother
\crefname{originalremark}{original remark}{original remarks}
\Crefname{originalremark}{Original remark}{Original remarks}
\makeatletter
\newlabel{remark-set-theoretic}{{0AIS}{source}{}{sourcepointer-11}{}}
\newlabel{remark-set-theoretic@cref}{{[originalremark][0][]0AIS}{[1][0][]source}{}{}{}}
\makeatother
\title{Fibre products of formal algebraic spaces preserve countable indexedness, classicity, weak adicness, adic* and local Noetherianness}
\author{The Stacks Project Authors}
\date{stacks-project:formal-spaces.tex:local-snapshot-20261012; GNU FDL 1.2 or later}
\begin{document}
\maketitle
\noindent\textit{Editorial note.} All delivered mathematical text is the author's own, copied from the pinned Stacks Project chapter file formal-spaces.tex; no statement or proof step was rewritten, reordered or removed.\par\medskip
\noindent\textit{Editorial note.} The chapter's own preamble.tex is shipped with the frozen source and defines the reference, slogan and history environments as comment environments, so material in those environments is discarded and citations inside them do not need a bibliography; the counted spans of this card contain no citation at all.\par\medskip
\noindent\textit{Editorial note.} The proof contains two of the author's own slips in wording, 'morphsms' for 'morphisms' and 'parapgrah' for 'paragraph'; both lie inside the counted span and are reproduced verbatim, as the delivered text is the author's.\par\medskip
\noindent\textit{Editorial note.} Parts (2), (3) and (4) are concluded by reference to separate parts of the same completed tensor product lemma, which is the author's own organisation; the version of that lemma used for each part is recorded as a prerequisite.\par\medskip
\noindent\textit{Editorial note.} The counted spans carry 12 navigation stubs; no bibliography citation occurs in the spans.\par\medskip
\noindent\textit{Editorial note.} The preamble span points at the chapter's own preamble.tex shipped with the frozen source, used verbatim as staged.\par\medskip
\noindent\textit{Allowed standard prerequisites (editorial).}\begin{itemize}
\item The definition of a formal algebraic space via coverings by affine formal algebraic spaces. Same chapter, Definition definition-formal-algebraic-space, invoked at proof lines 4817-4822. Supplies the coverings whose fibre products reduce the five statements to the affine case.
\item Affine formal algebraic spaces, and the fibre product of affines over a separated base being affine. Same chapter, Definition definition-affine-formal-algebraic-space at proof lines 4840-4842 and Lemma lemma-fibre-product-affines-over-separated at proof lines 4825-4827 and 4865-4867. Identifies the fibre product of the members of the coverings as affine, and later identifies the fibre product of formal spectra with the formal spectrum of the completed tensor product.
\item Countably indexed affine formal algebraic spaces as formal spectra, and morphisms between them as continuous ring maps. Same chapter, Lemma lemma-countably-indexed at proof lines 4858-4862, Lemma lemma-countable-affine-formal-algebraic-space at proof lines 4836-4837 and Lemma lemma-morphism-between-formal-spectra at proof lines 4862-4865. Converts the geometric statement into a statement about completed tensor products of topological rings.
\item The completed tensor product of weakly admissible topological rings and the properties it preserves. Same chapter, Lemma lemma-completed-tensor-product at proof lines 4867-4874 and again at proof lines 4876-4900, and Lemma lemma-completed-tensor-product-weakly-adic at proof lines 4897-4899. Supplies each of the four transfer statements that conclude parts (1) to (4).
\item That a morphism of schemes factors through a stage of a colimit presentation, and the reduction of an affine formal algebraic space. Same chapter, Lemma lemma-factor-through-thickening at proof lines 4844-4846, Lemma lemma-reduction-smooth at proof lines 4924-4927 and Example example-reduction-affine-formal-spectrum at proof lines 4949-4952. Underlies the cofinality construction and the translation of the finite type hypothesis in part (5).
\item Localness of being locally of finite type for morphisms of algebraic spaces. Morphisms of Spaces, Lemma spaces-morphisms-lemma-finite-type-local, invoked at proof lines 4909-4910. Moves the finite type hypothesis from the original morphism of reductions to the morphisms between the reductions of the members of the covering.
\end{itemize}\medskip
\ifdefined\hypertarget\hypertarget{sourcepointer-1}{}\fi
\noindent\textit{Source reference.} Original definition 0GB6: sources/stacks-formal-spaces/pinned-original-source.tex, source lines 2911--2911. Pointer outside the counted span, invoked once for context; the load-bearing argument of the counted statement is the author's proof reproduced in full in the proof span.\par\medskip
\ifdefined\hypertarget\hypertarget{sourcepointer-2}{}\fi
\noindent\textit{Source reference.} Original lemma 0GB4: sources/stacks-formal-spaces/pinned-original-source.tex, source lines 1187--1187. Pointer outside the counted span, invoked once for context; the load-bearing argument of the counted statement is the author's proof reproduced in full in the proof span.\par\medskip
\ifdefined\hypertarget\hypertarget{sourcepointer-3}{}\fi
\noindent\textit{Source reference.} Original lemma 0GXJ: sources/stacks-formal-spaces/pinned-original-source.tex, source lines 1756--1756. Pointer outside the counted span, invoked once for context; the load-bearing argument of the counted statement is the author's proof reproduced in full in the proof span.\par\medskip
\ifdefined\hypertarget\hypertarget{sourcepointer-4}{}\fi
\noindent\textit{Source reference.} Original lemma 0AII: sources/stacks-formal-spaces/pinned-original-source.tex, source lines 2423--2423. Pointer outside the counted span, invoked once for context; the load-bearing argument of the counted statement is the author's proof reproduced in full in the proof span.\par\medskip
\ifdefined\hypertarget\hypertarget{sourcepointer-5}{}\fi
\noindent\textit{Source reference.} Original lemma 0AN1: sources/stacks-formal-spaces/pinned-original-source.tex, source lines 2493--2493. Pointer outside the counted span, invoked once for context; the load-bearing argument of the counted statement is the author's proof reproduced in full in the proof span.\par\medskip
\ifdefined\hypertarget\hypertarget{sourcepointer-6}{}\fi
\noindent\textit{Source reference.} Original lemma 0AIA: sources/stacks-formal-spaces/pinned-original-source.tex, source lines 2080--2080. Pointer outside the counted span, invoked once for context; the load-bearing argument of the counted statement is the author's proof reproduced in full in the proof span.\par\medskip
\ifdefined\hypertarget\hypertarget{sourcepointer-7}{}\fi
\noindent\textit{Source reference.} Original lemma 0AN3: sources/stacks-formal-spaces/pinned-original-source.tex, source lines 3703--3703. Pointer outside the counted span, invoked once for context; the load-bearing argument of the counted statement is the author's proof reproduced in full in the proof span.\par\medskip
\ifdefined\hypertarget\hypertarget{sourcepointer-8}{}\fi
\noindent\textit{Source reference.} Original lemma 0AN0: sources/stacks-formal-spaces/pinned-original-source.tex, source lines 2355--2355. Pointer outside the counted span, invoked once for context; the load-bearing argument of the counted statement is the author's proof reproduced in full in the proof span.\par\medskip
\ifdefined\hypertarget\hypertarget{sourcepointer-9}{}\fi
\noindent\textit{Source reference.} Original lemma 0GB7: sources/stacks-formal-spaces/pinned-original-source.tex, source lines 2923--2923. Pointer outside the counted span, invoked once for context; the load-bearing argument of the counted statement is the author's proof reproduced in full in the proof span.\par\medskip
\ifdefined\hypertarget\hypertarget{sourcepointer-10}{}\fi
\noindent\textit{Source reference.} Original lemma 040Y: sources/stacks-chow/tags.txt, source lines 5273--5273. Pointer outside the counted span, invoked once for context; the load-bearing argument of the counted statement is the author's proof reproduced in full in the proof span.\par\medskip
\ifdefined\hypertarget\hypertarget{sourcepointer-11}{}\fi
\noindent\textit{Source reference.} Original remark 0AIS: sources/stacks-formal-spaces/pinned-original-source.tex, source lines 2769--2769. Pointer outside the counted span, invoked only from a delivered context span as a forward reference; the load-bearing argument of the counted statement is the author's proof reproduced in full in the proof span.\par\medskip
\begin{definition}
\label{definition-affine-formal-algebraic-space}
Let $S$ be a scheme. We say a sheaf $X$ on $(\Sch/S)_{fppf}$ is an
{\it affine formal algebraic space} if there exist
\begin{enumerate}
\item a directed set $\Lambda$,
\item a system $(X_\lambda, f_{\lambda \mu})$ over $\Lambda$
in $(\Sch/S)_{fppf}$ where
\begin{enumerate}
\item each $X_\lambda$ is affine,
\item each $f_{\lambda \mu} : X_\lambda \to X_\mu$ is a thickening,
\end{enumerate}
\end{enumerate}
such that
$$
X \cong \colim_{\lambda \in \Lambda} X_\lambda
$$
as fppf sheaves and $X$ satisfies a set theoretic condition
(see Remark \ref{remark-set-theoretic}). A
{\it morphism of affine formal algebraic spaces}
over $S$ is a map of sheaves.
\end{definition}

\begin{definition}
\label{definition-formal-algebraic-space}
Let $S$ be a scheme. We say a sheaf $X$ on $(\Sch/S)_{fppf}$ is a
{\it formal algebraic space} if there exist a family of maps
$\{X_i \to X\}_{i \in I}$ of sheaves such that
\begin{enumerate}
\item $X_i$ is an affine formal algebraic space,
\item $X_i \to X$ is representable by algebraic spaces and \'etale,
\item $\coprod X_i \to X$ is surjective as a map of sheaves
\end{enumerate}
and $X$ satisfies a set theoretic condition
(see Remark \ref{remark-set-theoretic}). A
{\it morphism of formal algebraic spaces}
over $S$ is a map of sheaves.
\end{definition}

\begin{lemma}
\label{lemma-types-fibre-products}
Let $S$ be a scheme. Let $X \to Y$ and $Z \to Y$ be
morphisms of formal algebraic space over $S$. Then
\begin{enumerate}
\item If $X$ and $Z$ are locally countably indexed, then $X \times_Y Z$
is locally countably indexed.
\item If $X$ and $Z$ are locally countably indexed and classical,
then $X \times_Y Z$ is locally countably indexed and classical.
\item If $X$ and $Z$ are weakly adic, then $X \times_Y Z$
is weakly adic.
\item If $X$ and $Z$ are locally adic*, then $X \times_Y Z$ is
locally adic*.
\item If $X$ and $Z$ are locally Noetherian and $X_{red} \to Y_{red}$
is locally of finite type, then $X \times_Y Z$ is locally Noetherian.
\end{enumerate}
\end{lemma}

\begin{proof}
Choose a covering $\{Y_j \to Y\}$ as in
Definition \ref{definition-formal-algebraic-space}.
For each $j$ choose a covering $\{X_{ji} \to Y_j \times_Y X\}$
as in Definition \ref{definition-formal-algebraic-space}.
For each $j$ choose a covering $\{Z_{jk} \to Y_j \times_Y Z\}$
as in Definition \ref{definition-formal-algebraic-space}.
Observe that $X_{ji} \times_{Y_j} Z_{jk}$ is an
affine formal algebraic space by
Lemma \ref{lemma-fibre-product-affines-over-separated}.
Hence
$$
\{X_{ji} \times_{Y_j} Z_{jk} \to X \times_Y Z\}
$$
is a covering as in Definition \ref{definition-formal-algebraic-space}.
Thus it suffices to prove (1), (2), (3), and (4) in case $X$, $Y$, and $Z$
are affine formal algebraic spaces.

\medskip\noindent
Assume $X$ and $Z$ are countably indexed. Say $X = \colim X_n$ and
$Z = \colim Z_m$ as in
Lemma \ref{lemma-countable-affine-formal-algebraic-space}.
Write $Y = \colim_{\lambda \in \Lambda} Y_\lambda$ as in
Definition \ref{definition-affine-formal-algebraic-space}.
For each $n$ and $m$ we can find $\lambda_{n, m} \in \Lambda$
such that $X_n \to Y$ and $Z_m \to Y$ factor through $Y_{\lambda_{n, m}}$
(for example see Lemma \ref{lemma-factor-through-thickening}).
Pick $\lambda_0 \in \Lambda$. By induction for $t \geq 1$
pick an element $\lambda_t \in \Lambda$ such that
$\lambda_t \geq \lambda_{n, m}$ for all $1 \leq n, m \leq t$
and $\lambda_t \geq \lambda_{t - 1}$. Set $Y' = \colim Y_{\lambda_t}$.
Then $Y' \to Y$ is a monomorphism such that $X \to Y$ and $Z \to Y$
factor through $Y'$. Hence we may replace $Y$ by $Y'$, i.e.,
we may assume that $Y$ is countably indexed.

\medskip\noindent
Assume $X$, $Y$, and $Z$ are countably indexed. By
Lemma \ref{lemma-countably-indexed} we can write
$X = \text{Spf}(A)$, $Y = \text{Spf}(B)$, $Z = \text{Spf}(C)$
for some weakly admissible topological rings $A$, $B$, and $C$.
The morphsms $X \to Y$ and $Z \to Y$ are given by
continuous ring maps $B \to A$ and $B \to C$, see
Lemma \ref{lemma-morphism-between-formal-spectra}.
By Lemma \ref{lemma-fibre-product-affines-over-separated}
we see that $X \times_Y Z = \text{Spf}(A \widehat{\otimes}_B C)$
and that $A \widehat{\otimes}_B C$ is a weakly admissible topological ring.
In particular, we see that $X \times_Y Z$ is countably indexed
by Lemma \ref{lemma-completed-tensor-product} part (3).
This proves (1).

\medskip\noindent
Proof of (2). In this case $X$ and $Z$ are countably indexed and
hence the arguments above show that $X \times_Y Z$ is the formal spectrum
of $A \widehat{\otimes}_B C$ where $A$ and $C$ are admissible.
Then $A \widehat{\otimes}_B C$ is admissible by 
Lemma \ref{lemma-completed-tensor-product} part (2).

\medskip\noindent
Proof of (3). As before we conclude that $X \times_Y Z$ is the formal spectrum
of $A \widehat{\otimes}_B C$ where $A$ and $C$ are weakly adic.
Then $A \widehat{\otimes}_B C$ is weakly adic by 
Lemma \ref{lemma-completed-tensor-product-weakly-adic}.

\medskip\noindent
Proof of (4). Arguing as above, this follows from
Lemma \ref{lemma-completed-tensor-product} part (4).

\medskip\noindent
Proof of (5). To deduce case (5) from
Lemma \ref{lemma-completed-tensor-product} part (5)
we need to show the hypotheses match.
Namely, with notation as in the first parapgrah of the proof,
if $X_{red} \to Y_{red}$ is locally of finite type, then
$(X_{ji})_{red} \to (Y_j)_{red}$ is locally of finite type. This follows
from Morphisms of Spaces, Lemma \ref{spaces-morphisms-lemma-finite-type-local}
and the fact that in the commutative diagram
$$
\xymatrix{
(X_{ji})_{red} \ar[d] \ar[r] & (Y_j)_{red} \ar[d] \\
X_{red} \ar[r] & Y_{red}
}
$$
the vertical morphisms are \'etale. Namely, we have
$(X_{ji})_{red} = X_{ij} \times_X X_{red}$ and
$(Y_j)_{red} = Y_j \times_Y Y_{red}$
by Lemma \ref{lemma-reduction-smooth}. Thus as above we reduce
to the case where $X$, $Y$, $Z$ are affine formal algebraic spaces,
$X$, $Z$ are Noetherian, and $X_{red} \to Y_{red}$ is of finite type.
Next, in the second paragraph of the proof we replaced $Y$ by $Y'$
but by construction $Y_{red} = Y'_{red}$, hence the finite type
assumption is preserved by this replacement. Then we see that
$X, Y, Z$ correspond to $A, B, C$ and $X \times_Y Z$ to
$A \widehat{\otimes}_B C$ with $A$, $C$ Noetherian adic.
Finally, taking the reduction corresponds to dividing by the
ideal of topologically nilpotent elements
(Example \ref{example-reduction-affine-formal-spectrum})
hence the fact that $X_{red} \to Y_{red}$ is of finite type
does indeed mean that $B/\mathfrak b \to A/\mathfrak a$ is of finite type
and the proof is complete.
\end{proof}

\end{document}
